\documentclass[conference]{IEEEtran}
\IEEEoverridecommandlockouts
% The preceding line is only needed to identify funding in the first footnote. If that is unneeded, please comment it out.
\usepackage{cite}
\usepackage{amsmath,amssymb,amsfonts}
\usepackage{algorithmic}
\usepackage{graphicx}
\usepackage{textcomp}
\usepackage{xcolor}
\usepackage[T1]{fontenc}
\def\BibTeX{{\rm B\kern-.05em{\sc i\kern-.025em b}\kern-.08em
    T\kern-.1667em\lower.7ex\hbox{E}\kern-.125emX}}
\begin{document}

\title{Conference Paper Title*\\
{\footnotesize \textsuperscript{*}Note: Sub-titles are not captured in Xplore and
should not be used}
\thanks{Identify applicable funding agency here. If none, delete this.}
}

\author{\IEEEauthorblockN{1\textsuperscript{st} Given Name Surname}
\IEEEauthorblockA{\textit{dept. name of organization (of Aff.)} \\
\textit{name of organization (of Aff.)}\\
City, Country \\
email address or ORCID}
\and
\IEEEauthorblockN{2\textsuperscript{nd} Given Name Surname}
\IEEEauthorblockA{\textit{dept. name of organization (of Aff.)} \\
\textit{name of organization (of Aff.)}\\
City, Country \\
email address or ORCID}
\and
\IEEEauthorblockN{3\textsuperscript{rd} Given Name Surname}
\IEEEauthorblockA{\textit{dept. name of organization (of Aff.)} \\
\textit{name of organization (of Aff.)}\\
City, Country \\
email address or ORCID}
\and
\IEEEauthorblockN{4\textsuperscript{th} Given Name Surname}
\IEEEauthorblockA{\textit{dept. name of organization (of Aff.)} \\
\textit{name of organization (of Aff.)}\\
City, Country \\
email address or ORCID}
\and
\IEEEauthorblockN{5\textsuperscript{th} Given Name Surname}
\IEEEauthorblockA{\textit{dept. name of organization (of Aff.)} \\
\textit{name of organization (of Aff.)}\\
City, Country \\
email address or ORCID}
\and
\IEEEauthorblockN{6\textsuperscript{th} Given Name Surname}
\IEEEauthorblockA{\textit{dept. name of organization (of Aff.)} \\
\textit{name of organization (of Aff.)}\\
City, Country \\
email address or ORCID}
}

\maketitle

\begin{abstract}
This document is a model and instructions for \LaTeX.
This and the IEEEtran.cls file define the components of your paper [title, text, heads, etc.]. *CRITICAL: Do Not Use Symbols, Special Characters, Footnotes, 
or Math in Paper Title or Abstract.
\end{abstract}

\begin{IEEEkeywords}
component, formatting, style, styling, insert
\end{IEEEkeywords}

\section{Introduction}
This document is a model and instructions for \LaTeX.
Please observe the conference page limits. 

\section{Revisão da Literatura}

\subsection{Qualidade perceptual e qualidade para a máquina: duas escalas
distintas}

A suposição implícita em boa parte dos pipelines de visão computacional é a de
que uma imagem melhor para o observador humano é também uma imagem melhor para
o classificador. A literatura revisada mostra que essa equivalência não se
sustenta. Dodge e Karam \cite{b1} foram os primeiros a formular o problema em
termos explícitos, ao distinguir a qualidade perceptual humana da qualidade do
ponto de vista de um sistema automático e ao argumentar que distorções
cotidianas de aquisição, transmissão e armazenamento são muito mais prováveis,
na prática, do que perturbações adversárias construídas deliberadamente.
Avaliando quatro arquiteturas pré-treinadas sobre o conjunto de validação do
ImageNet degradado por cinco famílias de distorção, os autores documentam uma
assimetria que reaparece de forma sistemática em toda a literatura posterior:
as redes são muito sensíveis a borramento e ruído e notavelmente resilientes a
compressão moderada e a alterações de contraste. Mais relevante para o presente
trabalho, registram que em desvios-padrão de ruído nos quais o desempenho médio
das redes cai abaixo de 20\%, as imagens permanecem facilmente reconhecíveis
por observadores humanos.

Karahan \textit{et al.} \cite{b2} replicam essa dissociação em reconhecimento
facial e a tornam ainda mais aguda. Aplicando cinco famílias de degradação
apenas ao conjunto de sondas do protocolo de identificação fechada do LFW,
observam que a acurácia \textit{rank}-1 de três CNNs pré-treinadas cai a menos
de 6\% sob níveis de ruído nos quais, segundo o relato dos autores, o olho
humano ainda identifica os sujeitos sem dificuldade; e, no sentido inverso,
verificam que as redes são surpreendentemente resilientes ao escurecimento, uma
distorção que prejudica o observador humano. As duas escalas de qualidade,
portanto, não apenas divergem em magnitude: divergem em ordenação.

Há também uma razão formal para esperar tensão entre aparência e fidelidade.
Blau e Michaeli \cite{b3} demonstram que qualidade perceptual e distorção são
objetivos mutuamente incompatíveis, com uma fronteira de compromisso
monotonicamente não crescente e convexa, resultado que vale para qualquer
medida de distorção --- inclusive distâncias definidas em espaços de
características de redes neurais. Convém registrar, contudo, que esse trabalho
não mede acurácia de classificação: o ``classificador'' do teorema é um
discriminador ótimo entre imagens reais e sintéticas, e a parte experimental
compara algoritmos de super-resolução por métricas de qualidade de imagem. A
ponte entre o compromisso percepção--distorção e o desempenho de
reconhecimento é, portanto, uma inferência que a presente revisão assume
explicitamente como hipótese, e não um resultado estabelecido. Sua utilidade é
delimitar a expectativa teórica: se otimizar a aparência necessariamente afasta
a imagem reconstruída da imagem original, e se o classificador depende de
propriedades dessa imagem original, então a degradação de desempenho tem uma
explicação anterior a qualquer análise espectral.

\subsection{Evidência empírica: transformações que melhoram a imagem degradam o
classificador}

O corpo empírico revisado converge para a conclusão de que as transformações
clássicas de entrada cobram um preço em acurácia, mesmo quando aplicadas a
imagens que não sofreram qualquer degradação prévia. A magnitude desse preço,
porém, varia por ordens de grandeza conforme a técnica, o dataset e o protocolo,
e é essa variação --- e não a existência do efeito --- que ainda carece de
explicação.

O borramento é o caso mais bem documentado. Vasiljevic \textit{et al.}
\cite{b4} avaliam a VGG-16 pré-treinada sobre versões borradas do conjunto de
validação do ImageNet e reportam queda monotônica da acurácia \textit{top}-5 de
90,88\% para 88,06\%, 81,48\%, 72,69\% e 60,97\% com desfoque de raio 2, 4, 6 e
8, chegando a 17,39\% com desvio gaussiano $\sigma = 8$. Um detalhe do seu
resultado merece destaque para a presente investigação: o borramento de
movimento de comprimento comparável é substancialmente menos danoso (83,32\%
contra 60,97\%), e os autores atribuem a diferença ao fato de a atenuação de
alta frequência ser unidirecional em um caso e radialmente uniforme no outro.
Isto é, a mesma ``quantidade'' de suavização produz danos distintos conforme a
geometria da supressão espectral. Roy \textit{et al.} \cite{b5} ampliam a
varredura para seis arquiteturas contra seis modelos de degradação e encontram
decaimento monotônico em todos os casos, com a observação adicional de que as
arquiteturas mais rasas resistem melhor do que as mais profundas.

O denoising --- a técnica cuja finalidade declarada é justamente melhorar a
imagem --- produz os resultados mais contraintuitivos. Hendrycks e Dietterich
\cite{b6} aplicam \textit{Non-Local Means} como etapa de pré-processamento
sobre o ImageNet-C e registram que o erro médio de corrupção sobe de 76,7\%
para 82,1\%: a restauração torna o classificador globalmente pior. Liao
\textit{et al.} \cite{b7} quantificam o custo diretamente sobre imagens não
atacadas, mostrando que um autoencoder \textit{denoising} padrão derruba a
acurácia do Inception-v3 de 76,7\% para 58,3\% --- 18,4 pontos percentuais
perdidos sem que houvesse qualquer degradação a remover. Paranhos da Costa
\textit{et al.} \cite{b8}, em um desenho de generalização cruzada completo,
observam que filtragem de mediana e \textit{Non-Local Means} recuperam parte do
desempenho perdido sobre entradas ruidosas mas jamais alcançam o \textit{baseline}
limpo, e atribuem explicitamente a lacuna residual à perda de detalhe e textura
causada pela suavização. Yim e Sohn \cite{b9} chegam a uma conclusão negativa
equivalente com uma Inception-v3: o pré-processamento isolado não produz
melhoria consistente e o modelo único alimentado com a imagem filtrada perde
desempenho sobre o dataset original.

A compressão e a quantização exibem um padrão qualitativamente diferente, de
platô seguido de queda. Dziugaite \textit{et al.} \cite{b10} fornecem o dado
primário, muito reutilizado na literatura posterior, de que a compressão JPEG
reduz a confiança média do classificador mesmo sobre imagens limpas. Prakash
\textit{et al.} \cite{b11} elevam essa observação a diagnóstico do campo,
afirmando que o maior obstáculo às defesas baseadas em transformação da entrada
é precisamente o fato de a transformação degradar a qualidade das imagens não
adversárias. Ainda assim, ambos os trabalhos e também \cite{b1} registram que a
compressão só se torna danosa abaixo de um limiar de qualidade relativamente
baixo --- um comportamento não linear que contrasta com a degradação monotônica
observada sob borramento.

Por fim, dois trabalhos recentes deslocam a evidência para pipelines de
reconstrução de estado da arte. Chen e Ma \cite{b12} documentam que o PyNET e o
LW-ISP, algoritmos de ISP inteligente projetados para melhorar a qualidade da
imagem, produzem desempenho de reconhecimento inferior ao do ISP tradicional da
câmera (37,2 contra 40,8 mAP no LOD, e resultados análogos em segmentação de
instâncias). Seu experimento mais valioso é de controle: fixando a qualidade de
reconstrução em 21,30 dB de PSNR, ainda obtêm 2,0 mAP de diferença conforme o
conteúdo espectral da reconstrução --- ou seja, o desempenho varia com a
distribuição de frequências mesmo quando a qualidade perceptual é mantida
constante. Hao \textit{et al.} \cite{b13}, em evidência secundária de detecção
de objetos, reforçam o mesmo descolamento ao observar que é possível comprimir
imagens em amplitude de forma substancial sem afetar a acurácia de detecção,
mesmo quando a qualidade visual já se deteriorou.

\subsection{A hipótese espectral: o que a CNN usa e o observador humano não vê}

A explicação que melhor articula esse conjunto de resultados sustenta que
classificadores convolucionais exploram, como sinal discriminativo, componentes
da imagem que o sistema visual humano descarta --- e que as técnicas de
aprimoramento perceptual, por construção, atuam exatamente sobre esses
componentes.

A formulação operacional mais antiga é a de Jo e Bengio \cite{b14}, que
propõem medir se as CNNs aprendem conceitos semânticos ou apenas regularidades
estatísticas superficiais. Aplicando máscaras radiais e aleatórias no domínio de
Fourier, calibradas visualmente para preservar a reconhecibilidade humana dos
objetos, os autores observam que o erro de uma \textit{Preact}-ResNet-200 sobe
de 5,22\% para 23,37\% no CIFAR-10. A conclusão é que a rede depende de energia
espectral de que o observador humano prescinde, e que sua boa generalização se
apoia em pistas superficiais compartilhadas entre treino e teste. Wang
\textit{et al.} \cite{b15} tornam a hipótese explícita e a ancoram em uma série
de experimentos de filtragem passa-baixa e passa-alta: enquanto humanos
dependem de componentes semânticos de baixa frequência, as CNNs minimizam a
perda explorando também sinais imperceptíveis de alta frequência. Mostram ainda
que heurísticas que aumentam a acurácia limpa, como \textit{Batch Normalization}
e \textit{Mix-up}, fazem-no incentivando a captura de alta frequência, e que
forçar a suavidade dos \textit{kernels} convolucionais --- que, pelo teorema da
convolução, equivale a impor um filtro passa-baixa --- derruba a acurácia em
dados limpos de 85,6\% para 63,1\%.

Geirhos \textit{et al.} \cite{b16} fornecem a verificação causal mais forte do
mecanismo, ainda que formulada no domínio espacial. Com 48.560 ensaios
psicofísicos sobre imagens de pistas conflitantes, estabelecem que CNNs
treinadas no ImageNet classificam por textura local enquanto humanos
classificam por forma global, e demonstram por intervenção sobre os dados de
treino que esse viés é induzido pelo dataset e não pela arquitetura. Se a
decisão da rede se apoia em textura, então qualquer transformação que suavize
\end{document}
textura remove o próprio sinal de decisão --- que é precisamente a explicação
oferecida de forma independente por \cite{b1}, \cite{b2} e \cite{b8}. Ilyas
\textit{et al.} \cite{b17} complementam o argumento pelo lado da utilidade:
características imperceptíveis ao humano são genuinamente preditivas, e removê-las
do conjunto de treino custa cerca de dez pontos de acurácia padrão no CIFAR-10
(85,40\% contra 95,25\%). É importante registrar, porém, que esse trabalho não
realiza qualquer análise espectral --- a ponte entre ``características
imperceptíveis'' e ``alta frequência'' é feita por terceiros \cite{b18}, e deve
ser tratada como tal.

A sensibilidade espectral tem ainda uma origem arquitetural. Tsuzuku e Sato
\cite{b19} derivam, a partir da linearização de camadas convolucionais com
\textit{padding} circular, que seus vetores singulares são combinações de poucas
funções de base de Fourier, e verificam empiricamente, em seis datasets e oito
arquiteturas, que existem ``pontos sensíveis'' bem definidos no espectro nos
quais uma perturbação de amplitude fixa altera a predição. O controle
metodológico decisivo é que um MLP treinado sobre os mesmos dados não exibe
esses pontos: a sensibilidade é atribuível à operação de convolução. Esse
resultado é relevante para a presente investigação porque o mecanismo é
independente da origem da alteração espectral --- vale igualmente para uma
perturbação adversária e para um filtro de pré-processamento.

\subsection{Da dependência espectral à previsão da degradação}

O passo seguinte --- estabelecer uma relação quantificável entre energia
espectral alterada e magnitude da queda de acurácia --- é justamente onde a
literatura se torna mais escassa e menos convergente.

A evidência mais próxima de uma relação dose--resposta explícita é a de Yin
\textit{et al.} \cite{b18}. Os autores caracterizam cada corrupção do CIFAR-10-C
e do ImageNet-C por uma métrica escalar, a fração de energia de alta frequência
do delta introduzido pela corrupção, e relacionam esse valor à variação de
acurácia por regressão linear, reportando inclinação e resíduo. Seus mapas de
calor de Fourier mostram que o modelo treinado naturalmente é sensível a
perturbações em quase todo o espectro, exceto nas frequências mais baixas, e que
aumento gaussiano e treino adversário deslocam essa dependência para as baixas
frequências, melhorando a robustez a corrupções de alta frequência e piorando-a
a corrupções de baixa frequência. Estabelecem também uma assimetria estrutural
que organiza vários resultados dispersos: o desempenho cai muito mais
rapidamente quando se remove informação de baixa frequência do que de alta,
porque imagens naturais concentram sua energia nas baixas frequências --- é mais
fácil para o modelo aprender a ignorar o que é escasso do que o que é
abundante.

Saikia \textit{et al.} \cite{b20} convertem essa ideia em um eixo ordinal,
ordenando as quinze corrupções do ImageNet-C por conteúdo de alta frequência e
mostrando um cruzamento sistemático: o modelo regularizado para robustez de alta
frequência é uniformemente melhor na metade superior do eixo e uniformemente
pior na inferior. Quantificam também o custo dessa escolha em acurácia limpa ---
o erro sobe de 11,8\% para 14,8\% e depois 16,0\% conforme o viés se intensifica
--- e registram uma justificativa que dialoga diretamente com a hipótese desta
revisão: preferem regularizar as ativações a filtrar a entrada porque o
\textit{denoising} destrói sinais de alta frequência importantes e pode
introduzir novos artefatos. Frank \textit{et al.} \cite{b21} oferecem a curva
dose--resposta mais limpa em um domínio específico, mostrando que a acurácia de
um classificador de imagens sintéticas cai de 98,24\% para 85,96\% e 84,20\% à
medida que o \textit{kernel} de \textit{anti-aliasing} cresce --- uma relação
direta entre intensidade de suavização espectral e desempenho, com a ressalva de
que, nesse domínio, o sinal discriminativo é de alta frequência por construção.

Contra a leitura simples dessa hipótese pesam, contudo, quatro resultados
importantes. Primeiro, Maiya \textit{et al.} \cite{b22}, treinando trinta e seis
modelos por dataset com mascaramento probabilístico de bandas via DCT, mostram
que a banda crítica depende do dataset e não é universalmente a alta frequência:
remover 90\% da banda 0--15 custa cerca de 2\% de acurácia no CIFAR-10, 10\% no
TinyImageNet e 35\% no ImageNet, enquanto remoções altíssimas na banda 48--63
têm efeito praticamente nulo fora do CIFAR-10. Segundo, Tsuzuku e Sato
\cite{b19} encontram que as bandas mais danosas são as baixas e médias, não as
altas. Terceiro, Ortiz-Jiménez \textit{et al.} \cite{b23} demonstram que a
magnitude da degradação depende menos da energia removida do que de \emph{onde}
ela é removida no protocolo: filtrar passa-baixa apenas o conjunto de teste
derruba a acurácia do CIFAR-10 de cerca de 93\% para 27,45\%, enquanto treinar
com os mesmos dados filtrados custa cerca de três pontos. A explicação que
oferecem --- as fronteiras de decisão só existem nas direções em que o
classificador foi treinado com características discriminativas --- separa com
nitidez duas causas que a literatura frequentemente confunde: perda de
informação e descasamento de domínio. Quarto, o próprio \cite{b18} reporta o
contraexemplo em que sua hipótese falha, ao mostrar que treinar com ruído de
estatísticas de Fourier idênticas às da corrupção de névoa piora o desempenho
sob névoa; e Shao \textit{et al.} \cite{b24} observam que a filtragem tradicional
sequer produz o deslocamento espectral limpo que se lhe atribui, porque apenas
reduz a energia das bandas indesejadas em vez de eliminá-la, deixando informação
residual explorável.

\subsection{Mecanismos concorrentes e complementares}

A explicação espectral não é a única disponível, e os trabalhos que propõem
alternativas ajudam a delimitar o que ela efetivamente explica.

O mais direto é o descasamento de distribuição. Vasiljevic \textit{et al.}
\cite{b4} mostram que o ajuste fino com imagens borradas quase fecha a lacuna de
desempenho, o que permite descartar tanto a deficiência arquitetural quanto a
ausência de informação discriminativa na imagem borrada: o sinal semântico
sobrevive, mas passa a ser mal representado. A análise de ativações que realizam
localiza o problema com precisão, ao mostrar que a invariância ao borramento só
é adquirida nas camadas profundas, enquanto a disparidade nas camadas iniciais
permanece alta. Paranhos da Costa \textit{et al.} \cite{b8} chegam à mesma
conclusão por outra via, ao observar que modelos treinados sobre imagens
restauradas generalizam mal para imagens ruidosas --- o pré-processamento cria
um novo domínio em vez de restaurar o original.

Um segundo mecanismo opera no nível do filtro individual. Borkar e Karam
\cite{b25} ordenam os filtros convolucionais por suscetibilidade e mostram que
corrigir apenas metade deles em uma única camada já recupera parte substancial
do desempenho, indicando que o dano não se distribui uniformemente. Sua
evidência mais valiosa para esta revisão é qualitativa mas contundente: os
filtros mais atingidos pelo borramento gaussiano são precisamente os seletivos
em frequência e orientação, isto é, os detectores de borda e contorno, enquanto
o ruído branco atinge também os \textit{blobs} específicos de cor. Há, portanto,
uma correspondência previsível entre a natureza espectral da transformação e a
família de filtros que ela desativa.

Um terceiro mecanismo, proposto por Liao \textit{et al.} \cite{b7}, é o efeito de
amplificação de erro. Medindo a perturbação relativa camada a camada, os autores
mostram que o resíduo deixado por um \textit{denoiser} otimizado em espaço de
pixels, embora pequeno na entrada, é progressivamente amplificado ao longo da
hierarquia até atingir, no topo, magnitude próxima à de uma perturbação
adversária. Esse resultado explica por que menor erro de reconstrução não implica
melhor classificação --- constatação replicada por \cite{b26}, que registra que a
aplicação excessiva de \textit{denoising} remove detalhes cruciais e produz
borramento mesmo quando as métricas de qualidade parecem favoráveis. Dodge e
Karam \cite{b1} já haviam observado o fenômeno da propagação em termos
qualitativos, notando que o borramento altera pouco as respostas da primeira
camada convolucional e muito as da última, ao passo que o ruído é captado
imediatamente pelos filtros iniciais: duas vias distintas de dano na mesma
hierarquia.

Convém registrar, finalmente, que parte da literatura de defesa adversária opera
sobre a premissa inversa --- a de que a perturbação relevante é de alta
frequência e que descartá-la é, portanto, seguro. Entezari e Papalexakis
\cite{b27} enunciam essa premissa de forma particularmente limpa ao justificar a
aproximação de baixo posto como mecanismo de defesa. É exatamente essa premissa
que os trabalhos com instrumentação espectral direta \cite{b18}, \cite{b22}
contestam, e que Maiya \textit{et al.} \cite{b22} descrevem como consolidada por
viés de confirmação metodológico, uma vez que não é possível testar o cenário
inverso --- bloquear as baixas frequências --- sem destruir a informação
associada ao rótulo.

\subsection{Condições em que o pré-processamento melhora}

A literatura também delimita, com razoável precisão, as condições sob as quais a
intervenção sobre a entrada não degrada ou até beneficia o classificador, e essas
condições são informativas por contraste.

A fronteira mais nítida é traçada por Hendrycks e Dietterich \cite{b6}, que
avaliam duas técnicas clássicas sob protocolo idêntico e obtêm sinais opostos:
enquanto o \textit{Non-Local Means} eleva o erro médio de corrupção de 76,7\%
para 82,1\%, a equalização adaptativa de histograma (CLAHE) o reduz para 74,5\%
e melhora inclusive o erro limpo, de 23,87\% para 23,55\%. A distinção que os
autores propõem é qualitativa e não de grau: o CLAHE redistribui intensidades,
ao passo que o \textit{denoising} descarta detalhe fino. O critério que separa
melhoria de degradação, portanto, não é alterar a imagem, mas remover
informação.

As demais condições de melhoria registradas na literatura confirmam esse
critério. O \textit{denoising} beneficia a classificação quando a entrada está
genuinamente degradada, o método é casado ao tipo de ruído e a dose é controlada
\cite{b26}, \cite{b8}; deixa de beneficiar quando o método é agressivo demais e
dilui as bordas do objeto \cite{b9}. A imagem filtrada agrega desempenho quando
usada como canal adicional ao lado da original, e não quando a substitui
\cite{b9} --- o mesmo padrão que Geirhos \textit{et al.} \cite{b16} encontram em
outro nível: treinar somente no Stylized-ImageNet custa dezesseis pontos de
\textit{top}-1, mas treinar em SIN e IN conjuntamente supera a linha de base
tanto em acurácia limpa quanto em robustez. O borramento moderado funciona como
aumento de dados e chega a melhorar a acurácia sobre imagens nítidas
\cite{b4}, e a restauração explícita da entrada por \textit{deblurring}, ainda
que recupere boa parte do desempenho, é inferior e muito mais custosa do que
adaptar o modelo à degradação. A manipulação espectral usada como augmentação de
treino, e não como filtro de teste, melhora simultaneamente o desempenho limpo e
o degradado \cite{b28}, \cite{b29}. E a transformação de entrada é tolerável
quando descarta seletivamente: limiares adaptativos por imagem em base
\textit{wavelet} reduzem o custo sobre imagens limpas a menos de dois pontos
percentuais, contra os custos muito maiores de tabelas de quantização fixas e
agnósticas ao conteúdo \cite{b11}. Há, por fim, a condição trivial e
esclarecedora identificada por Ortiz-Jiménez \textit{et al.} \cite{b23}: quando
o dataset simplesmente não contém informação discriminativa na banda removida,
como no MNIST sob filtragem passa-baixa de largura $14\times14$, a filtragem não
custa nada.

O padrão que emerge desse conjunto é coerente. O pré-processamento é benigno
quando redistribui em vez de remover, quando é aplicado simetricamente a treino e
teste, quando preserva a imagem original como canal complementar, ou quando atua
sobre bandas que o dataset não usa. É danoso quando remove, assimetricamente,
energia de bandas nas quais o classificador aprendeu a decidir.

\subsection{Lacunas e posicionamento deste trabalho}

A síntese dos trabalhos revisados permite identificar quatro lacunas que
motivam a presente investigação.

A primeira é de escopo de técnica. O corpo empírico concentra-se quase
inteiramente em degradações --- borramento, ruído, compressão, redução de
resolução --- e em \textit{denoising}, que é a única técnica de finalidade
restauradora avaliada com regularidade \cite{b6}, \cite{b7}, \cite{b8},
\cite{b9}, \cite{b26}. Técnicas cuja finalidade declarada é o aprimoramento
perceptual --- equalização de histograma, aguçamento e super-resolução --- são
raras: a equalização aparece em um único experimento de configuração única
\cite{b6}, o aguçamento não é avaliado em nenhum dos estudos revisados, e a
super-resolução é tratada exclusivamente sob a ótica da qualidade de imagem
\cite{b3}, sem medição de desempenho de classificação. A pergunta central deste
trabalho --- se transformações \emph{perceptualmente benéficas} degradam
classificadores --- permanece, portanto, testada apenas em uma fração do espaço
de técnicas que ela abrange.

A segunda é de variável independente. Mesmo os estudos que varrem múltiplos
níveis de intensidade parametrizam a dose pelo controle da técnica --- tamanho
de \textit{kernel}, desvio padrão, fator de quantização, raio de filtragem
\cite{b1}, \cite{b4}, \cite{b5}, \cite{b28} --- e não pela energia espectral
efetivamente removida. A única exceção é \cite{b18}, que constrói uma métrica
escalar de fração de energia de alta frequência e a relaciona à variação de
acurácia por regressão. Não existe, na literatura revisada, uma curva
dose--resposta entre energia espectral removida por banda e queda de acurácia,
construída sobre técnicas clássicas de aprimoramento perceptual e classificadores
pré-treinados em modo de inferência. Duas transformações com o mesmo PSNR podem
distribuir a perda de energia de formas inteiramente distintas \cite{b13}, e é
precisamente essa distinção que as métricas de distorção usuais não capturam.

A terceira é de confundimento entre perda de informação e descasamento de
domínio. O achado de \cite{b23} --- três pontos de custo ao filtrar o treino
contra sessenta e cinco ao filtrar apenas o teste --- mostra que boa parte da
degradação atribuída à remoção de sinal é, de fato, deslocamento de
distribuição. A maioria dos estudos de degradação não separa essas duas
contribuições, o que torna as magnitudes reportadas difíceis de comparar entre
protocolos. Avaliar classificadores pré-treinados em modo estrito de inferência,
sem qualquer reajuste, isola deliberadamente o cenário de interesse prático ---
aquele em que o pipeline de pré-processamento é inserido a montante de um modelo
que não será retreinado --- mas exige que a interpretação dos resultados
reconheça essa ambiguidade em vez de ignorá-la.

A quarta é de fragilidade da própria hipótese espectral. A banda crítica depende
do dataset \cite{b22}; as bandas mais danosas nem sempre são as altas
\cite{b19}; casar as estatísticas de Fourier de uma corrupção não basta para
prever seu efeito \cite{b18}; a dependência espectral do modelo é herdada da
distribuição espectral dos dados de treino, e não é uma propriedade fixa da
arquitetura \cite{b16}; e o deslocamento da dependência espectral pode
dissociar-se do desempenho \cite{b24}. Uma investigação que pretenda explicar a
degradação pela alteração da distribuição de energia espectral precisa,
portanto, tratar essa relação como hipótese a ser testada com medição direta ---
e não como premissa herdada --- e precisa estar preparada para quantificar
quanto da variância observada ela de fato explica.

É nessa interseção que o presente trabalho se situa: aplicar técnicas clássicas
de pré-processamento e aprimoramento perceptual a imagens naturais, medir a
variação de acurácia de classificadores CNN pré-treinados em modo de inferência,
caracterizar cada transformação pela energia espectral que efetivamente remove ou
introduz por banda de frequência, e verificar em que medida essa caracterização
prevê a magnitude da degradação observada.

% =====================================================================
% Referências --- acrescentar ao bloco \begin{thebibliography} do template
% =====================================================================

\begin{thebibliography}{00}

\bibitem{b1} S. Dodge and L. Karam, ``Understanding how image quality affects deep neural networks,'' arXiv preprint arXiv:1604.04004, 2016.

\bibitem{b2} S. Karahan, M. K. Yıldırım, K. Kırtaç, F. Ş. Rende, G. Bütün, and H. K. Ekenel, ``How image degradations affect deep CNN-based face recognition?,'' arXiv preprint arXiv:1608.05246, 2016.

\bibitem{b3} Y. Blau and T. Michaeli, ``The perception-distortion tradeoff,'' in Proc. IEEE Conf. Computer Vision and Pattern Recognition (CVPR), 2018, pp. 6228--6237.

\bibitem{b4} I. Vasiljevic, A. Chakrabarti, and G. Shakhnarovich, ``Examining the impact of blur on recognition by convolutional networks,'' arXiv preprint arXiv:1611.05760, 2017.

\bibitem{b5} P. Roy, S. Ghosh, S. Bhattacharya, and U. Pal, ``Effects of degradations on deep neural network architectures,'' arXiv preprint arXiv:1807.10108, 2018.

\bibitem{b6} D. Hendrycks and T. Dietterich, ``Benchmarking neural network robustness to common corruptions and perturbations,'' in Proc. Int. Conf. Learning Representations (ICLR), 2019.

\bibitem{b7} F. Liao, M. Liang, Y. Dong, T. Pang, X. Hu, and J. Zhu, ``Defense against adversarial attacks using high-level representation guided denoiser,'' arXiv preprint arXiv:1712.02976, 2018.

\bibitem{b8} G. B. Paranhos da Costa, W. A. Contato, T. S. Nazaré, J. E. S. Batista Neto, and M. Ponti, ``An empirical study on the effects of different types of noise in image classification tasks,'' arXiv preprint arXiv:1609.02781, 2016.

\bibitem{b9} J. Yim and K.-A. Sohn, ``Enhancing the performance of convolutional neural networks on quality degraded datasets,'' arXiv preprint arXiv:1710.06805, 2017.

\bibitem{b10} G. K. Dziugaite, Z. Ghahramani, and D. M. Roy, ``A study of the effect of JPG compression on adversarial images,'' arXiv preprint arXiv:1608.00853, 2016.

\bibitem{b11} A. Prakash, N. Moran, S. Garber, A. DiLillo, and J. Storer, ``Deflecting adversarial attacks with pixel deflection,'' arXiv preprint arXiv:1801.08926, 2018.

\bibitem{b12} H. Chen and K. Ma, ``Enhancing vision: Harmonizing frequency for imaging quality and perception accuracy,'' in Proc. IEEE Int. Conf. Acoustics, Speech and Signal Processing (ICASSP), 2025.

\bibitem{b13} Y. Hao, Y. Pei, Y. Lyu, Z. Yuan, J.-R. Rizzo, Y. Wang, and Y. Fang, ``Understanding the impact of image quality and distance of objects to object detection performance,'' arXiv preprint arXiv:2209.08237, 2022.

\bibitem{b14} J. Jo and Y. Bengio, ``Measuring the tendency of CNNs to learn surface statistical regularities,'' arXiv preprint arXiv:1711.11561, 2017.

\bibitem{b15} H. Wang, X. Wu, Z. Huang, and E. P. Xing, ``High-frequency component helps explain the generalization of convolutional neural networks,'' in Proc. IEEE/CVF Conf. Computer Vision and Pattern Recognition (CVPR), 2020.

\bibitem{b16} R. Geirhos, P. Rubisch, C. Michaelis, M. Bethge, F. A. Wichmann, and W. Brendel, ``ImageNet-trained CNNs are biased towards texture; increasing shape bias improves accuracy and robustness,'' in Proc. Int. Conf. Learning Representations (ICLR), 2019.

\bibitem{b17} A. Ilyas, S. Santurkar, D. Tsipras, L. Engstrom, B. Tran, and A. Mądry, ``Adversarial examples are not bugs, they are features,'' in Advances in Neural Information Processing Systems (NeurIPS), 2019.

\bibitem{b18} D. Yin, R. Gontijo Lopes, J. Shlens, E. D. Cubuk, and J. Gilmer, ``A Fourier perspective on model robustness in computer vision,'' in Advances in Neural Information Processing Systems (NeurIPS), 2019.

\bibitem{b19} Y. Tsuzuku and I. Sato, ``On the structural sensitivity of deep convolutional networks to the directions of Fourier basis functions,'' in Proc. IEEE/CVF Conf. Computer Vision and Pattern Recognition (CVPR), 2019.

\bibitem{b20} T. Saikia, C. Schmid, and T. Brox, ``Improving robustness against common corruptions with frequency biased models,'' arXiv preprint arXiv:2103.16241, 2021.

\bibitem{b21} J. Frank, T. Eisenhofer, L. Schönherr, A. Fischer, D. Kolossa, and T. Holz, ``Leveraging frequency analysis for deep fake image recognition,'' in Proc. Int. Conf. Machine Learning (ICML), 2020.

\bibitem{b22} S. R. Maiya, M. Ehrlich, V. Agarwal, S.-N. Lim, T. Goldstein, and A. Shrivastava, ``A frequency perspective of adversarial robustness,'' arXiv preprint arXiv:2111.00861, 2021.

\bibitem{b23} G. Ortiz-Jiménez, A. Modas, S.-M. Moosavi-Dezfooli, and P. Frossard, ``Hold me tight! Influence of discriminative features on deep network boundaries,'' in Advances in Neural Information Processing Systems (NeurIPS), 2020.

\bibitem{b24} Z. Shao, Y. Ren, X. Wang, L. Isik, and D. Beck, ``Dissociating spatial frequency reliance from adversarial robustness advantages in neurally guided deep convolutional neural networks,'' manuscrito não publicado.

\bibitem{b25} T. S. Borkar and L. J. Karam, ``DeepCorrect: Correcting DNN models against image distortions,'' arXiv preprint arXiv:1705.02406, 2019.

\bibitem{b26} M. Hami and M. J. JameBozorg, ``Assessing the impact of CNN auto encoder-based image denoising on image classification tasks,'' arXiv preprint arXiv:2404.10664, 2024.

\bibitem{b27} N. Entezari and E. E. Papalexakis, ``TensorShield: Tensor-based defense against adversarial attacks on images,'' arXiv preprint arXiv:2002.10252, 2020.

\bibitem{b28} J. Wu, ``FCR: Robust image classification method with frequency domain consistency regularization,'' in Proc. 6th Int. Symp. Computer Technology and Information Science (ISCTIS), IEEE, 2026.

\bibitem{b29} ``Improving corruption robustness with random erasing in the frequency domain,'' 2023.

\end{thebibliography}
